%==============================================================================
\section{Networking API}
\label{sec:api}
%==============================================================================
This section specifies the interface between the madGLP runtime and the networking layer. The interface is independent of the underlying communication medium.
%------------------------------------------------------------------------------
\subsection{Identity}
\label{sec:identity}
%------------------------------------------------------------------------------
Each agent has an Ed25519 key pair. The 32-byte public key serves as the agent's globally unique identity, used for addressing, authentication, and signing. Identity is persistent across sessions.
\begin{longtable}{L{4cm} L{4.5cm} L{5.5cm}}
\toprule
\rowcolor{headerblue}
\textcolor{white}{\textbf{Function}} &
\textcolor{white}{\textbf{Signature}} &
\textcolor{white}{\textbf{Description}} \\
\midrule
\endhead
Generate identity & \func{putIdentity()} & Generate and persist an Ed25519 key pair. \\
\rowcolor{rowgray}
Retrieve identity & \func{getIdentity()} $\to$ \code{PubKey} & Return the agent's public key. \\
\bottomrule
\end{longtable}
%------------------------------------------------------------------------------
\subsection{Peer Discovery and Connection}
\label{sec:discovery}
%------------------------------------------------------------------------------
The networking layer notifies the runtime when peers become reachable or unreachable, regardless of the underlying medium.

When two agents cannot discover each other via BLE proximity, they may exchange identities out-of-band using \func{generatePeerLink()} and \func{consumePeerLink(uri)}.  A peer link is a signed, time-limited bearer invitation: possession of a valid unused link authorizes one first-contact attempt by the holder.
\begin{longtable}{L{4cm} L{5cm} L{5cm}}
\toprule
\rowcolor{headerblue}
\textcolor{white}{\textbf{Function}} &
\textcolor{white}{\textbf{Signature}} &
\textcolor{white}{\textbf{Description}} \\
\midrule
\endhead
Scan for peers & \func{getPeers()} $\to$ \code{List<Peer>} & Discover currently reachable peers. \\
\rowcolor{rowgray}
Peer discovered & \func{onPeerDiscovered(cb)} & Callback when a new peer is detected. The callback receives the peer's public key and nickname.  \\

Peer connected & \func{onPeerConnected(cb)} & Callback when a peer becomes reachable for communication. \\
\rowcolor{rowgray}
Peer disconnected & \func{onPeerDisconnected(cb)} & Callback when a peer is no longer reachable. \\
Check reachability & \func{isPeerReachable(pk)} $\to$ \code{bool} & Whether the given peer is currently reachable via any medium. \\
\rowcolor{rowgray}
Connectivity status & \func{onConnectivityStatus(cb)} & Callback when the networking layer detects a change in the agent's public IP address.  Triggers the rendezvous protocol (Section~\ref{sec:rendezvous}).  Fired on startup and on address change. \\
Generate peer link & \func{generatePeerLink()} $\to$ \code{URI} & Generate a signed, one-time, RV-mediated peer-link invitation for out-of-band sharing.  Full semantics are specified in Appendix~\ref{appendix:ip-cold-call-invites}. \\
\rowcolor{rowgray}
Consume peer link & \func{consumePeerLink(uri)} $\to$ \code{PubKey} & Verify and redeem an incoming peer-link invitation through a listed rendezvous server, register the inviter's identity, and attempt first contact.  Returns the inviter's public key; see Appendix~\ref{appendix:ip-cold-call-invites}. \\
\bottomrule
\end{longtable}
%------------------------------------------------------------------------------
\subsection{Cold Calls and Trusted Environment Levels}
\label{sec:cold-call-trust}
%------------------------------------------------------------------------------
A \emph{cold call} is a first-contact connection attempt between two agents that do not already share an active grassroots networking communication channel. Cold calls may be initiated by a nearby BLE encounter or by an explicit deep-link shared out-of-band, requiring IP connectivity (Section~\ref{sec:ip-cold-call}).  In the IP case, the inviting agent need not know the redeemer's public key before the link is opened.  A cold call is implicitly a friendship request: accepting the cold call establishes friendship.

Agents may enforce one of two cold-call trust levels, affecting BLE encounters:
\begin{itemize}
\item \textbf{Open:} The agent surfaces unsolicited cold calls from nearby BLE peers as friendship requests for the user to accept or decline.
\item \textbf{Closed:} The agent ignores unsolicited cold calls from nearby BLE encounters without surfacing them.
\end{itemize}
Explicit cold-calls through deep-links shared out-of-band are not affected by the different trust levels: the link is itself an explicit invitation by its creator, so the resulting friendship request is always surfaced for the user to accept or decline regardless of the configured trust level.  The invite link includes authorization: it contains a temporary, one-time 256-bit token generated by its creator that is consumed upon usage; see Appendix~\ref{appendix:ip-cold-call-invites} for the whole invite flow.

In the \textbf{Open} trust level, an agent responds to any nearby Grassroots BLE encounter by completing the ANNOUNCE exchange, thereby revealing its full public key and allowing first contact. In the \textbf{Closed} trust level, an agent does not complete ANNOUNCE with an unknown nearby device: it uses the derived UUID
to recognize already-known peers, and ignores unmatched BLE encounters. 

%------------------------------------------------------------------------------
\subsection{Point-to-Point Communication}
\label{sec:p2p}
%------------------------------------------------------------------------------
The networking layer provides reliable delivery of arbitrary byte payloads between peers.
\begin{longtable}{L{4cm} L{5.5cm} L{4.5cm}}
\toprule
\rowcolor{headerblue}
\textcolor{white}{\textbf{Function}} &
\textcolor{white}{\textbf{Signature}} &
\textcolor{white}{\textbf{Description}} \\
\midrule
\endhead
Send payload & \func{send(pubKey, payload)} & Send a byte payload to the peer identified by \code{pubKey}. \\
\rowcolor{rowgray}
Receive payload & \func{onReceive(cb)} & Register callback for incoming payloads. The callback receives the sender's public key and the byte payload. \\
\bottomrule
\end{longtable}
%------------------------------------------------------------------------------
\subsection{Networking Assumptions}
\label{sec:assumptions}
%------------------------------------------------------------------------------
The networking layer must satisfy the following assumptions for madGLP correctness:
\textbf{Fair message delivery:} Every message passed to \func{send()} is eventually delivered to the recipient, assuming the recipient eventually becomes reachable. Messages to temporarily-unreachable peers are queued and delivered when connectivity is restored.
\textbf{Atomic message delivery:} Each message is delivered as a complete unit. If a message is fragmented across multiple packets, the networking layer reassembles them before delivery.
\textbf{Integrity:} Messages are not corrupted in transit. The networking layer provides encryption and integrity checks.
\textbf{Authentication:} The sender's identity (public key) is authenticated. A received message from public key $pk$ was indeed sent by the agent holding the corresponding private key.
